\documentclass[xcolor=table,aspectratio=169,fleqn]{beamer}
\usepackage{beamerthemesplit}
\usepackage{wrapfig}
\usetheme{SPbGU}
\usepackage{pdfpages}
\usepackage{amsmath}
\usepackage{mathtools}
\usepackage{cmap}
\usepackage[T2A]{fontenc}
\usepackage[utf8]{inputenc}
\usepackage[english]{babel}
\usepackage{indentfirst}
\usepackage{newtxmath}
\usepackage{tikz}
\usepackage{multirow}
\usepackage[noend]{algpseudocode}
\usepackage{algorithm}
\usepackage{algorithmicx}
\usepackage{fancyvrb}
\usepackage{hyperref}
\usepackage{nicematrix}
\definecolor{links}{HTML}{2A1B81}
\hypersetup{colorlinks,linkcolor=,urlcolor=links}
\usetikzlibrary{calc}
\usetikzlibrary{shapes, backgrounds}
\usetikzlibrary{arrows,automata}
\usetikzlibrary{positioning}
\usetikzlibrary{fit}
\usetikzlibrary{shapes.callouts}
\usetikzlibrary{shapes.misc}
\usepackage{xparse}
\usepackage{fontawesome}

\usepackage{etoolbox,refcount}
\usepackage{multicol}

\usepackage{tabularx}
\newcolumntype{Y}{>{\raggedleft\arraybackslash}X}

\renewcommand{\thealgorithm}{}

\newtheorem{mytheorem}{Theorem}
\renewcommand{\thealgorithm}{}

\newcommand{\tikzmark}[1]{\tikz[overlay,remember picture] \node (#1) {};}
\def\Put(#1,#2)#3{\leavevmode\makebox(0,0){\put(#1,#2){#3}}}

\newcommand{\ltz}{$< 1$}

\tikzset{
    state/.style={
           rectangle,
           rounded corners,
           draw=black, very thick,
           minimum height=1.5em,
           inner sep=2pt,
           text centered,
           },
}

\tikzset{
    invisible/.style={opacity=0,text opacity=0},
    visible on/.style={alt=#1{}{invisible}},
    alt/.code args={<#1>#2#3}{%
      \alt<#1>{\pgfkeysalso{#2}}{\pgfkeysalso{#3}} % \pgfkeysalso doesn't change the path
    },
}

\tikzset{cross/.style={cross out, draw=black, minimum size=2*(#1-\pgflinewidth), inner sep=0pt, outer sep=0pt, ultra thick},
%default radius will be 1pt. 
cross/.default={1pt}}

\NewDocumentCommand{\mycallout}{r<> O{opacity=0.8,text opacity=1} m m m}{%
\tikz[remember picture, overlay]\node[align=center, fill=cyan!20, text width=#5cm,
#2,visible on=<#1>, rounded corners,
draw,rectangle callout,anchor=pointer,callout relative pointer={(290:0.5cm)}]
at (#3) {#4};
}

\NewDocumentCommand{\mycalloutR}{r<> O{opacity=0.8,text opacity=1} m m m}{%
\tikz[remember picture, overlay]\node[align=center, fill=cyan!20, text width=#5cm,
#2,visible on=<#1>, rounded corners,
draw,rectangle callout,anchor=pointer,callout relative pointer={(30:0.8cm)}]
at (#3) {#4};
}

\newcommand\colR{\cellcolor{red!20}}
\newcommand\colB{\cellcolor{blue!20}}
\newcommand\colG{\cellcolor{green!20}}
\newcommand\colO{\cellcolor{orange!20}}
\definecolor{Gray}{gray}{0.8}

%callout relative pointer={(230:0.5cm)}]

\newcounter{countitems}
\newcounter{nextitemizecount}
\newcommand{\setupcountitems}{%
  \stepcounter{nextitemizecount}%
  \setcounter{countitems}{0}%
  \preto\item{\stepcounter{countitems}}%
}
\makeatletter
\newcommand{\computecountitems}{%
  \edef\@currentlabel{\number\c@countitems}%
  \label{countitems@\number\numexpr\value{nextitemizecount}-1\relax}%
}
\newcommand{\nextitemizecount}{%
  \getrefnumber{countitems@\number\c@nextitemizecount}%
}
\newcommand{\previtemizecount}{%
  \getrefnumber{countitems@\number\numexpr\value{nextitemizecount}-1\relax}%
}
\makeatother    
\newenvironment{AutoMultiColItemize}{%
\ifnumcomp{\nextitemizecount}{>}{3}{\begin{multicols}{2}}{}%
\setupcountitems\begin{itemize}}%
{\end{itemize}%
\unskip\computecountitems\ifnumcomp{\previtemizecount}{>}{3}{\end{multicols}}{}}


\beamertemplatenavigationsymbolsempty


\title[Динамический анализ графов]{Динамический анализ графов}
\institute[СПбГУ]{
Санкт-Петербургский Государственный Университет
}

\author[Семён Григорьев]{Семён Григорьев}

\date{\today}

\begin{document}
{
\begin{frame}[fragile]
  \titlepage
\end{frame}
}

\begin{frame}[fragile]
  \frametitle{Немного ссылок}
  \begin{itemize}    
    \item Uri Zwick: \url{https://www.math.tau.ac.il/~zwick/}
    \item Giuseppe F. Italiano: \url{https://scholar.google.com/citations?hl=ru&user=6jiwt-UAAAAJ&view_op=list_works&sortby=pubdate}
    \item differential-dataflow: \url{https://github.com/TimelyDataflow/differential-dataflow}
    \item DDlog: \url{https://github.com/vmware-archive/differential-datalog}
    \item FlowLog: \url{https://github.com/flowlog-rs/flowlog}
    \item Viatra: \url{https://eclipse.dev/viatra/documentation/query-language.html}
    \pause
    \item \href{https://www.mdpi.com/2076-3417/15/11/6003}{Recent Advances in Efficient Dynamic Graph Processing}
    \item \href{https://arxiv.org/abs/2404.18211}{A survey of dynamic graph neural networks}
    \item \href{https://ieeexplore.ieee.org/abstract/document/11202740}{A Comprehensive Survey of Dynamic Graph Neural Networks: Models, Frameworks, Benchmarks, Experiments and Challenges}
    \item \href{https://bonndoc.ulb.uni-bonn.de/xmlui/handle/20.500.11811/10104}{Temporal Graph Algorithms}
    \item \href{https://www.sciencedirect.com/science/article/pii/S0167739X25003917}{ContraMST: A unified framework for dynamic MST maintenance}
    \item \href{https://dl.acm.org/doi/10.1145/3318464.3389733}{Regular Path Query Evaluation on Streaming Graphs}
    \item \href{https://people.cs.umass.edu/~immerman/pub/dynfo.pdf}{Dyn-FO: A Parallel, Dynamic Complexity Class}
  \end{itemize}
\end{frame}

\begin{frame}[fragile]
  \frametitle{Рандомизированные алгоритмы}
  \begin{itemize}    
    \item \textbf{Monte Carlo}: ошибается с некоторой (небольшой) вероятностью, то гарантии на время работы
    \item \textbf{Las Vegas}:точный ответ, но с некоторой вероятностью может работать медленно (на фоне <<хорошего среднего случая>>) или даже не останавливаться
  \end{itemize}
\end{frame}

%\begin{frame}[fragile]
%  \frametitle{Decremental reachability\footnote{Из перзентации <<Improved dynamic reachability algorithms for directed graphs>>, Liam Roditty и Uri Zwick}}
%  \begin{center}
%  \begin{tabular}{|c|c|c|c|c|}
%    \hline
%    Graphs & Graphs & Query time & Total\footnote{$m$ --- количество изменений} update time & Authors \\
%    \hline
%    \hline
%    DAGs & Deterministic & 1 & $mn$ & Italiano ’88 \\
%    General & Monte Carlo & $\frac{n}{\log n}$ & $m n \log^2 n $& Henzinger King ’95 \\
%    General & Deterministic & 1 & $m^2$ & La Poutré, van Leeuwen ’87 \\
%    General & Deterministic & 1 & $n^3$ & Demetrescu, Italiano ’00\\
%    General & Monte Carlo & 1 & $mn^{\frac{4}{3}}$ & Baswana, Hariharan, Sen ’02 \\
%    General & Las Vegas & 1 & $mn$ & Roditty, Zwick \\
%    \hline
%  \end{tabular}
%\end{center}
%\end{frame}
%
%\begin{frame}[fragile]
%  \frametitle{Fully dynamic reachability\footnote{Из перзентации <<Improved dynamic reachability algorithms for directed graphs>>, Liam Roditty и Uri Zwick}}
%  \begin{center}
%  \begin{tabular}{|c|c|c|c|c|}
%    \hline
%    Graphs & Graphs & Query time & Amortized\footnote{<<среднее>> время обработки одного изменения} update time & Authors \\
%    \hline
%    \hline
%    General & Deterministic & 1 & $n^2 \log n$ & King ’99 \\
%    General & Deterministic & 1 & $n^2 n $& Demetrescu, Italiano ’00 \\
%    General & Deterministic & 1 & $n^2$ & Roditty ’03 \\
%    General & Deterministic & 1 & $n^3$ & Demetrescu, Italiano ’00\\
%    General & Monte Carlo & 1 & $mn^{\frac{4}{3}}$ & Baswana, Hariharan, Sen ’02 \\
%    General & Las Vegas & 1 & $mn$ & Roditty, Zwick \\
%    \hline
%  \end{tabular}
%\end{center}
%\end{frame}

\end{document}
